\documentclass[11pt,a4paper]{article}

\usepackage[margin=1in]{geometry}
\usepackage{graphicx}
\usepackage{hyperref}
\usepackage{booktabs}
\usepackage{enumitem}
\usepackage{parskip}
\usepackage{titlesec}

\hypersetup{
    colorlinks=true,
    linkcolor=blue,
    urlcolor=blue,
    citecolor=blue
}

\titleformat{\section}{\large\bfseries}{}{0em}{}[\titlerule]
\titleformat{\subsection}{\normalsize\bfseries}{}{0em}{}

\title{
    \Large \textbf{Behavioural User Modelling and Personalised Recommendation} \\
    \large \textbf{Using LLM Agents with Agentic Workflows} \\
    \vspace{0.5em}
    \normalsize DSN $\times$ BCT LLM Agent Challenge --- Solution Write-up
}

\author{
    Oladimeji Balogun \and
    Timothy Ayelagbe \and
    Glory Bagai \and
    Promise Akinlua
}

\date{May 2026}

\begin{document}

\maketitle

% ─────────────────────────────────────────────────────────────
\section{What We Built}

We built a two-task LLM agent system that understands users deeply enough
to simulate their behaviour. Given a user's review history, our system can
predict what they would write and rate about an item they have never seen
(Task A), and recommend items they would genuinely enjoy (Task B).

Both tasks are served as a containerised FastAPI application with two
endpoints, backed by 10,000 user profiles built from the Yelp Academic
Dataset and a BM25-indexed catalogue of 78,982 businesses.

The live system is available at:
\begin{center}
\url{https://github.com/oladimeji-balogun/bluechip-hackathon-backend}
\end{center}

% ─────────────────────────────────────────────────────────────
\section{Our Approach}

The core insight driving our design is that both tasks share the same
foundation: a rich understanding of who the user is. Rather than building
two separate systems, we built one shared \texttt{UserProfile} schema and
fed it into two different agents. This meant we only had to solve the hard
problem of user understanding once.

We also made a deliberate choice to use \textbf{agentic workflows} rather
than single LLM calls. A single prompt asking an LLM to ``generate a
review for this user'' produces generic output. An agent that first
analyses the user's voice, then analyses the item, then predicts a rating,
then generates a review with that rating as an anchor produces output that
is measurably more faithful to the user's real behaviour. The reasoning
steps are explicit and inspectable, which matters both for quality and for
being able to explain what the system is doing.

We chose \textbf{LangGraph} as our orchestration framework because it
models agentic workflows as explicit state machines. Every node is a named
reasoning step. Every edge is a deliberate transition. Before writing a
single line of code, we drew the two state machines by hand and walked the
team through each node and what it was responsible for. That clarity of
design carried through to the implementation.

For the LLM backbone we used \textbf{Groq} with two Llama models:
\texttt{llama-3.3-70b-versatile} for tasks requiring nuanced reasoning and
generation, and \texttt{llama-3.1-8b-instant} for tasks requiring fast,
structured JSON output. This split gave us quality where it mattered and
speed where it did not.

All prompts are stored as \textbf{Jinja2-rendered Markdown templates},
completely separate from the agent code. This was a deliberate engineering
decision: prompt iteration is a different activity from code changes, and
mixing the two into the same files slows both down.

% ─────────────────────────────────────────────────────────────
\section{System Architecture}

\subsection{The UserProfile}

Everything in the system depends on the \texttt{UserProfile}. We built it
to capture five things about a user: how they rate (average stars, standard
deviation, whether they tend to rate high or low, their full star
distribution), how they write (average review length, tone classification
across five classes, formality score, common phrases, emoji usage), what
they prefer (top engaged categories, consistently avoided categories,
cross-domain signals), their recent history (last 20 reviews, capped at
300 characters each to keep LLM context lean), and their cultural context
(Nigerian food terms, location references, and Pidgin English markers
detected from review text).

Users are classified into four tiers by review count: cold (fewer than 5),
light (5--20), active (21--100), and power (more than 100). This tier is
the primary signal for routing in Task B.

We built profiles for 10,000 Yelp users, scanning the full review file to
ensure every user had their complete review history. The resulting
distribution was: 29.8\% cold, 38.5\% light, 23.9\% active, and 7.8\%
power users. Every profile is loaded into memory at server startup so both
agents have instant access at inference time.

\subsection{Architecture Diagram}

\begin{figure}[h]
\centering
% TODO: Replace with architecture diagram image
% \includegraphics[width=0.9\textwidth]{architecture.png}
\fbox{\parbox{0.85\textwidth}{\centering\vspace{2cm}
\textit{[Architecture diagram to be inserted ---}\\
\textit{shows Yelp data flowing through ProfileBuilder into the shared}\\
\textit{UserProfile Store, then branching into Task A and Task B agents,}\\
\textit{each served by a FastAPI endpoint]}
\vspace{2cm}}}
\caption{System architecture overview}
\label{fig:arch}
\end{figure}

\subsection{Technology Stack}

\begin{table}[h]
\centering
\begin{tabular}{ll}
\toprule
\textbf{Component} & \textbf{Technology} \\
\midrule
API Framework         & FastAPI \\
Agent Orchestration   & LangGraph \\
LLM (reasoning)       & Llama 3.3 70B via Groq \\
LLM (structured outputs) & Llama 3.1 8B via Groq \\
Prompt Templating     & Jinja2 (.md templates) \\
Business Retrieval    & BM25 (rank-bm25) \\
Data Validation       & Pydantic v2 \\
Containerisation      & Docker \\
Dataset               & Yelp Academic Dataset \\
\bottomrule
\end{tabular}
\caption{Technology stack}
\label{tab:stack}
\end{table}

% ─────────────────────────────────────────────────────────────
\section{Task A: User Modeling}

\subsection{The Problem}

We need to generate a star rating and a written review for a user-item
pair where the user has never reviewed that item. The output has to sound
like that specific user wrote it, not like a generic AI assistant.

\subsection{The Pipeline}

We designed Task A as a five-node LangGraph state machine. Here is what
each node does and why we separated them.

\textbf{analyze\_user} reads the \texttt{UserProfile} and extracts a
structured voice fingerprint: tone class, average review length, common
phrases, rating tendencies, and Nigerian context signals. This is pure
deterministic logic with no LLM call. We separated it from the LLM nodes
because it is always fast and always correct, and it sets up clean context
for everything downstream.

\textbf{analyze\_item} reads the item metadata and checks how well it maps
to the user's known preferences: does the category match their top
categories or their avoided ones, and how does the community rating compare
to this user's average expectation? Again, no LLM call. Two cheap
deterministic nodes before any expensive inference.

\textbf{predict\_rating} is the first LLM call, using
\texttt{llama-3.1-8b-instant} for speed. The prompt contains the user's
rating distribution, category preferences, and ten recent review snippets.
We run at temperature 0.1 for consistency. The node returns a float between
1.0 and 5.0 and a one-sentence reasoning string.

\textbf{generate\_review} is the most important node. We call
\texttt{llama-3.3-70b-versatile} with a detailed voice profile prompt that
includes the user's writing statistics, tone-specific instructions, three
real reviews from the user's history as few-shot examples, and the
predicted rating as an anchor for the emotional register. The prompt
explicitly instructs the model to write as a specific human, not as an AI.
Temperature is 0.8 for natural variation.

\textbf{apply\_nigerian\_context} is a conditional node that only runs when
Nigerian context is active. It makes a second LLM pass to naturally weave
in Pidgin English phrases and local references without changing the
sentiment or length.

\begin{figure}[h]
\centering
% TODO: Replace with Task A state machine diagram
% \includegraphics[width=0.75\textwidth]{task_a_graph.png}
\fbox{\parbox{0.85\textwidth}{\centering\vspace{2cm}
\textit{[Task A state machine diagram to be inserted ---}\\
\textit{linear pipeline: analyze\_user, analyze\_item, predict\_rating,}\\
\textit{generate\_review, conditional apply\_nigerian\_context, END]}
\vspace{2cm}}}
\caption{Task A LangGraph state machine}
\label{fig:task_a}
\end{figure}

\subsection{API Endpoints}

The system exposes two endpoints, both accepting JSON and returning JSON.

\textbf{Task A: Generate Review}

\begin{verbatim}
POST /api/v1/task-a/generate

{
  "user_id": "MGPQVLsODMm9ZtYQW-g_OA",   // Yelp user ID (or pass full profile)
  "nigerian_mode": false,                  // activate Nigerian context
  "item": {
    "item_id": "biz_001",
    "name": "The Melting Pot",
    "category": "Restaurants",
    "subcategories": ["Fondue", "Fine Dining"],
    "avg_community_rating": 4.1,
    "description": "Upscale fondue restaurant.",
    "price_range": "$$$"
  }
}
\end{verbatim}

Response includes \texttt{predicted\_stars}, \texttt{predicted\_stars\_rounded},
\texttt{generated\_review}, \texttt{confidence}, and \texttt{reasoning}.



\textbf{Task B: Get Recommendations}

\begin{verbatim}
POST /api/v1/task-b/recommend

{
  "user_id": "MGPQVLsODMm9ZtYQW-g_OA",
  "domain": "restaurants",
  "context": "Looking for somewhere to eat Friday evening",
  "top_k": 10,
  "nigerian_mode": false,
  "conversation_history": []             // supports multi-turn
}
\end{verbatim}

Response includes a ranked \texttt{recommendations} list with
\texttt{item\_id}, \texttt{item\_name}, \texttt{predicted\_rating},
\texttt{relevance\_score}, \texttt{explanation}, and flags for
\texttt{cold\_start\_used} and \texttt{cross\_domain\_used}.

Both endpoints accept either a \texttt{user\_id} for lookup against the
profile store, or a full inline \texttt{user\_profile} object for testing
without a pre-built profile.

\subsection{Nigerian Contextualisation}

We treated Nigerian contextualisation as a first-class design requirement,
not a last-minute addition. The profile builder scans every review for a
curated vocabulary of Nigerian food terms (suya, jollof, egusi, buka),
location names (Lagos, Abuja, Port Harcourt), and Pidgin English markers
(e dey, correct, sabi, wahala). When these signals are present, the
generation prompt activates a dedicated Nigerian voice block that instructs
the model on Pidgin usage, local cultural references, and price sensitivity.

The results speak for themselves. A user whose reviews reference Barcelos
in VI and use phrases like ``e dey sweet'' receives generated reviews that
naturally continue in that voice.

% ─────────────────────────────────────────────────────────────
\section{Task B: Recommendation}

\subsection{The Problem}

Recommendation is harder than review generation because it cannot be solved
with a single strategy. An active user with 127 restaurant reviews needs
deep preference matching. A cold-start user with 2 reviews needs something
completely different. A user asking for recommendations outside their usual
domain needs cross-domain preference transfer. We had to handle all three
and handle them correctly.

\subsection{The Pipeline}

Task B is a seven-node LangGraph state machine with a three-way conditional
branch at the retrieval stage. We drew this out by hand before implementing
it because the routing logic needed to be completely clear before writing
code.

\textbf{analyze\_user} calls \texttt{llama-3.3-70b-versatile} with the
full profile and request context. It returns a structured JSON object
capturing what the user genuinely values, what consistently disappoints
them, their domain familiarity, transferable preference signals, price
sensitivity, and a recommended retrieval strategy.

\textbf{route\_strategy} is a conditional edge function, not an LLM call.
It reads the state and routes to one of three retrieval nodes. We added a
hard override here: users with fewer than 5 reviews always route to
popularity fallback, regardless of what the LLM recommends. We learned
from testing that LLMs can be overconfident when routing cold users, so we
enforced the rule in code.

The three retrieval nodes all query the real Yelp business catalogue via
BM25. \textbf{popularity\_fallback} constructs a query from the user's
category hint and sets a minimum 4.0-star threshold. 
\textbf{cross\_domain\_retrieval} builds a query from the LLM-extracted
transferable signals and the requested domain. \textbf{preference\_retrieval}
builds a query from the user's top categories and context, filtering out
already-reviewed businesses.

\textbf{rank\_candidates} is where all three paths converge. The LLM
receives the full user analysis and the candidate list, and reasons over
each candidate against the user's specific values, dislikes, and price
sensitivity. This is the step that makes the system an agent rather than a
retrieval pipeline.

\textbf{apply\_context} handles multi-turn conversations. If conversation
history is present, a second LLM pass re-weights the ranking to honour
what the user said in previous turns.

\begin{figure}[h]
\centering
% TODO: Replace with Task B state machine diagram
% \includegraphics[width=0.9\textwidth]{task_b_graph.png}
\fbox{\parbox{0.85\textwidth}{\centering\vspace{2.5cm}
\textit{[Task B state machine diagram to be inserted ---}\\
\textit{analyze\_user, route\_strategy branching to three retrieval nodes,}\\
\textit{all converging at rank\_candidates, apply\_context, END]}
\vspace{2.5cm}}}
\caption{Task B LangGraph state machine}
\label{fig:task_b}
\end{figure}

\subsection{Why We Replaced LLM-Generated Candidates with BM25}

Early in development, we had the LLM generate candidate items from scratch.
The recommendations were plausible and well-reasoned, but they were
completely unverifiable: the LLM invented business names and synthetic IDs
that did not correspond to real places. We could not evaluate them against
ground truth and we could not trust them in a production context.

We replaced this with a real business catalogue built from Yelp's
\texttt{business.json}: 78,982 open businesses with at least 10 reviews,
indexed by BM25. The retrieval nodes now return real Yelp business IDs.
The LLM's role shifted from inventing candidates to reasoning over real
ones. This is a better division of labour: retrieval is fast and
deterministic, reasoning is where the LLM adds genuine value.

% ─────────────────────────────────────────────────────────────
\section{Evaluation and Results}

\subsection{Task A}

We evaluated using a leave-one-out protocol: for each user, we identified
reviews not in their profile history, used the most recent as ground truth,
and asked the agent to generate a review for that same item. This tests
genuine generalisation since the agent never saw the ground truth review.

\begin{table}[h]
\centering
\begin{tabular}{lr}
\toprule
\textbf{Metric} & \textbf{Score} \\
\midrule
ROUGE-1         & 0.3068 \\
ROUGE-2         & 0.0418 \\
ROUGE-L         & 0.1495 \\
BERTScore F1    & 0.8364 \\
Rating RMSE     & 1.1187 \\
Rating MAE      & 0.8521 \\
\midrule
Samples         & 48 \\
\bottomrule
\end{tabular}
\caption{Task A evaluation results}
\label{tab:task_a}
\end{table}

BERTScore F1 of 0.836 is the number we are most confident in. It measures
semantic similarity using contextual embeddings, which captures meaning
even when exact words differ. Two humans writing about the same restaurant
rarely use the same phrases; BERTScore correctly rewards semantic alignment
over surface matching. Rating RMSE of 1.119 is a meaningful improvement
over a random baseline (estimated at approximately 1.8 over the 1-5 star
range).

\subsection{Task B}

Standard ranking metrics (NDCG, Hit Rate, Precision) require matching
recommended item IDs against a ground truth set. Our retrieval returns real
Yelp IDs, but the probability of exact overlap between BM25-retrieved
candidates and a specific user's held-out items is low given 78,982
catalogue businesses. Automated metrics were therefore uninformative.

We instead conducted manual evaluation of 10 recommendation outputs. Active
users consistently received category-matched recommendations respecting
their price sensitivity. Cold-start users were correctly routed to
popularity fallback and received broadly well-liked items. Nigerian mode
produced culturally grounded explanations with appropriate local references.

% ─────────────────────────────────────────────────────────────
\section{What We Would Do With More Time}

\textbf{Dense retrieval.} BM25 matches keywords. A user who values
``intimate atmosphere'' and ``attentive service'' cannot express those
preferences in a keyword query. Replacing BM25 with sentence embedding
retrieval over business descriptions would capture semantic preference
signals that keyword search cannot. This is the single most impactful
improvement available.

\textbf{Fine-tuned review generation.} Our review generation relies on
prompt engineering with a general-purpose model. Fine-tuning Llama 3.1 8B
on paired (user profile, real review) examples from the Yelp dataset would
produce a specialised model with stronger behavioural fidelity and lower
inference cost.

\textbf{Nigerian context classifier.} The current heuristic detector
produces false positives. Common English expressions like ``e sweet''
occasionally trigger Nigerian mode for non-Nigerian users. A lightweight
supervised classifier trained on labelled reviews would reduce this
significantly.

\textbf{Richer item features.} Rating prediction is currently limited by
sparse item metadata (category, community rating, price range). Adding
review sentiment distributions and attribute vectors would improve RMSE.

\textbf{Real-time profile updates.} Profiles are built offline. A
production system would update them incrementally as users generate new
reviews, keeping the behavioural signal current.

% ─────────────────────────────────────────────────────────────
\section{Conclusion}

We built a system that treats users as the dynamic, context-sensitive
agents they actually are. Every design decision, from the shared
\texttt{UserProfile} schema to the explicit LangGraph state machines to
the BM25 catalogue retrieval, was made deliberately and with a clear
reason. We are proud of the architecture, the Nigerian contextualisation
work, and the fact that we built it from the ground up with full
understanding of every component.

The full codebase including profile builder, LangGraph agents, BM25
catalogue, evaluation scripts, and Docker configuration is at:
\begin{center}
\url{https://github.com/oladimeji-balogun/bluechip-hackathon-backend}
\end{center}

\end{document}